\documentclass[11pt,a4paper]{article}
\usepackage[T1]{fontenc}
\usepackage{lmodern}
\usepackage[margin=1in]{geometry}
\usepackage{amsmath,amssymb,amsthm,booktabs,array}
\usepackage[colorlinks=true,allcolors=blue]{hyperref}
\usepackage{xurl}
\hypersetup{pdftitle={Proof of Notik's Singer-ruler conjecture with sparse certificates},
pdfauthor={Haobo Ma; Wenlin Zhang},
pdfsubject={Largest-gap cuts of cyclic perfect difference sets and unique real chamber minima; CC BY 4.0},
pdfkeywords={Golomb ruler, Singer difference set, chamber, dual certificate}}
\setlength{\emergencystretch}{1em}
\newtheorem{theorem}{Theorem}
\newtheorem{lemma}[theorem]{Lemma}
\newtheorem{corollary}[theorem]{Corollary}
\newtheorem{proposition}[theorem]{Proposition}
\newcommand{\R}{\mathbb R}
\newcommand{\Z}{\mathbb Z}
\title{Proof of Notik's Singer-ruler conjecture\\with sparse certificates}
\author{Haobo Ma\\\small ChronoAI Pte Ltd, Singapore
\and Wenlin Zhang\\\small National University of Singapore\\\small ChronoAI Pte Ltd, Singapore}
\date{21 November 2026}
\begin{document}
\maketitle
\begin{abstract}
We prove Notik's conjecture in Section~7, page~7 of
\emph{Golomb rulers with real marks} (Zenodo version~2,
31~October 2026, DOI 10.5281/zenodo.23844781;
arXiv:2611.\allowbreak22656v1):
for every prime power, every Singer ruler is the unique shortest
point of its real, unit-separated distance-order chamber.
More generally, cutting any cyclic perfect difference set at its unique largest
gap gives this property. Every positive integer smaller than the removed
gap occurs as an ordinary distance; the chamber inequalities therefore
bound each retained adjacent gap below by its integer value.
We give a positive integer dual certificate supported on exactly one row
per retained gap; its matrix is unit triangular. The same argument gives
sharp stability bounds. Explicit largest-gap examples answer Notik's
analogous Bose and Ruzsa questions negatively.
\end{abstract}

\section{The question and the chamber}

Notik~\cite[Section~7, p.~7]{notik} states:
\begin{quote}
``For every prime power $m$, every Singer ruler is the only shortest point
of its chamber.''
\end{quote}
Here a Singer ruler means a cyclic Singer perfect difference set, possibly
multiplied by any unit, cut at a largest cyclic gap and translated to begin
at zero. Notik gives exact certificates for $24\,450$ distinct such rulers
of orders at most $103$. The argument below works for every order and
its unique largest-gap cut. Its conclusion concerns one chamber,
not the optimum over all real rulers.
Notik's earlier subset-sum work~\cite[Theorems~3.2, 4.3 and~4.4]{subset} uses
positive integer weights and triangular tight relations for the
Conway--Guy and Bohman families. Here the relations are obtained directly
from the ordinary distances immediately below the retained gaps.
The length formula $L=N-G$ for the shortest translated cut is
classical; see Dimitromanolakis~\cite[Section~5.7.1, pp.~48--49]{dimitro}
and the comment in OEIS A401165~\cite{oeisgap}.
The chamber problem allows all real marks with the same labelled
distance order, rather than just translations of a cyclic set.

Let $k\ge1$ and let $x=(x_0,\ldots,x_k)$ be an integer Golomb ruler:
$0=x_0<x_1<\cdots<x_k$, with all positive labelled distances
$\delta_{ij}(x)=x_j-x_i$ ($0\le i<j\le k$) distinct.
Its normalized unit-separated chamber $C(x)\subset\R^{k+1}$ consists of
the vectors $y$ satisfying
\begin{align}
 y_0&=0, & \delta_{ij}(y)&\ge1 &&(i<j),\label{eq:positive}\\
 \delta_{ab}(y)-\delta_{ij}(y)&\ge1
 &&\text{whenever }\delta_{ij}(x)<\delta_{ab}(x).\label{eq:order}
\end{align}
The objective is the length $y_k$. These are the chambers of
\cite[Section~2]{notik}, written in marks rather than consecutive gaps.
The underlying ordering of distances is the combinatorial equivalence
used by Beck, Bogart and Pham~\cite[Section~1]{beck}; the unit separation
in \eqref{eq:positive}--\eqref{eq:order} fixes the scale.
The integer vector $x$ itself belongs to $C(x)$, because distinct positive
integer distances are separated by at least one.

\begin{lemma}[A complete prefix of distances]\label{lem:prefix}
Put $a_i=x_i-x_{i-1}$ for $1\le i\le k$, and
$A=\max_i a_i$. If every integer $1,\ldots,A$ occurs as a positive
distance of $x$, then every $y\in C(x)$ satisfies
\[
 y_i-y_{i-1}\ge a_i\qquad(1\le i\le k).
\]
Consequently $x$ is the unique point of minimum length in $C(x)$.
\end{lemma}
\begin{proof}
For $1\le d\le A$, choose the labelled pair $(p_d,q_d)$ whose
reference distance is $d$. Define
$D_d(y)=y_{q_d}-y_{p_d}$.
The chamber gives $D_1(y)\ge1$ and
$D_d(y)-D_{d-1}(y)\ge1$ for $2\le d\le A$.
Adding yields $D_d(y)\ge d$.
The distance $a_i$ belongs to the pair $(i-1,i)$, so uniqueness of
reference distances gives $D_{a_i}(y)=y_i-y_{i-1}$.
Summation gives $y_k\ge\sum_i a_i=x_k$.
Equality forces equality in every adjacent-gap bound, since every
gap excess is nonnegative. With $y_0=x_0=0$, this gives $y=x$.
\end{proof}

\begin{corollary}[Stability and weighted objectives]\label{cor:stability}
Under the hypotheses of Lemma~\ref{lem:prefix}, put
$\varepsilon=y_k-x_k$ for $y\in C(x)$. Then
\[
 0\le y_i-x_i\le\varepsilon\quad(0\le i\le k),\qquad
 \sum_{i=1}^k\left|(y_i-y_{i-1})-a_i\right|=\varepsilon.
\]
Every nonnegative weighted sum of the gaps is minimized at $x$;
if all weights are positive, the minimizer is unique.
The coefficient $1$ in the coordinate bound is sharp.
\end{corollary}
\begin{proof}
The gap excesses are nonnegative and sum to $\varepsilon$.
Each mark excess is a partial sum of them. The weighted assertion
follows term by term, with equality forcing all gap excesses to vanish
when the weights are positive. The bound is an equality at $i=k$;
positive excess occurs, for example, at any dilation $y=tx$, $t>1$.
\end{proof}

\section{Largest-gap cuts}

A cyclic perfect difference set in $\Z/N\Z$, with $N>1$, is a subset $D$ such that
each nonzero residue has exactly one representation $b-a$ by an ordered
pair of distinct elements $a,b\in D$.
Its cyclic gaps are distinct: each is a nonzero residue represented by
its consecutive ordered pair, so equal gaps would contradict the perfect
difference property. Thus the largest cyclic gap $G$ is unique.
Cut there, and translate and sort the
remaining marks to obtain
\[
 0=x_0<x_1<\cdots<x_k=L<N,\qquad N=L+G.
\]
The retained adjacent gaps are strictly smaller than $G$.

\begin{theorem}\label{thm:main}
Every largest-gap cut of a cyclic perfect difference set is the unique
shortest point of its normalized unit-separated real chamber.
In fact every retained adjacent gap is minimized at that cut.
\end{theorem}
\begin{proof}
First, two equal ordinary positive distances would represent the same
nonzero residue modulo $N$, contradicting uniqueness of the ordered
pair. Thus $x$ is an integer Golomb ruler and $x\in C(x)$.

Let $1\le d<G$. The perfect difference property supplies a pair
$(i,j)$ such that $x_j-x_i\equiv d\pmod N$.
If $j<i$, its residue is $N-(x_i-x_j)\ge N-L=G$, a contradiction.
Hence $i<j$, and $0<x_j-x_i<N$ implies $x_j-x_i=d$.
Therefore every positive integer below $G$ is an ordinary distance.

Let $A$ be the largest retained gap. We have $A<G$.
Every $1\le d\le A$ therefore occurs by the preceding argument.
Lemma~\ref{lem:prefix} now proves the assertions.
\end{proof}

\begin{corollary}[Notik's conjecture]\label{cor:singer}
For every prime power $q$, every unit multiple of every Singer perfect
difference set modulo $q^2+q+1$, cut at its unique largest cyclic gap, is the
unique shortest point of its normalized real chamber.
\end{corollary}
\begin{proof}
Singer's construction~\cite{singer} gives a perfect difference set
with $q+1$ elements modulo $q^2+q+1$.
Multiplication by a unit permutes the nonzero difference residues and
preserves uniqueness of the representing pairs. Translation preserves
all differences. Theorem~\ref{thm:main} applies to each largest-gap cut.
For completeness, a finite-field construction establishing the needed
perfect difference property is given in the appendix.
\end{proof}

\section{An explicit dual certificate}

The same argument supplies the structured certificate requested in
\cite[Sections~7 and~10]{notik}.
Use gap coordinates $g=(g_1,\ldots,g_k)$, with
$g_i=y_i-y_{i-1}$. A labelled distance is its interval sum.
Let $v_d\in\{0,1\}^k$ be the interval indicator for the reference
distance $d$, for $1\le d\le A$, and set $v_0=0$.
For $d=1,\ldots,A$ put
\[
 r_d=v_d-v_{d-1},\qquad
 w_d=\#\{i:a_i\ge d\}.
\]
The rows $r_d\cdot g\ge1$ are valid chamber inequalities and are tight
at the reference gaps $a=(a_1,\ldots,a_k)$.

\begin{proposition}\label{prop:dual}
These rows have positive integer weights and satisfy
\[
 \sum_{d=1}^A w_dr_d=(1,\ldots,1),\qquad
 \sum_{d=1}^A w_d=\sum_{i=1}^k a_i=L.
\]
Their span is $\R^k$. Thus they certify both the minimum $L$ and
the uniqueness of the minimizer.
\end{proposition}
\begin{proof}
Every $w_d>0$, since some $a_i=A$.
The unique interval of distance $a_i$ is the single gap $i$, so
$v_{a_i}=e_i$, the $i$th coordinate vector. Telescoping gives
$\sum_{d=1}^{a_i}r_d=e_i$.
Adding these identities over $i$ gives the first formula; counting
the pairs $(i,d)$ with $1\le d\le a_i$ gives the second.
The coordinate-vector identities also prove the span assertion.
If $\sum_i g_i=L$, then
$\sum_d w_d(r_d\cdot g-1)=0$; positivity forces every row to be
tight. Their span then forces $g=a$.
\end{proof}

For example, $\{0,1,3\}\subset\Z/7\Z$ has cut gap $G=4$ and
retained gaps $(1,2)$. Here $A=2$, $v_1=(1,0)$,
$v_2=(0,1)$, $r_1=(1,0)$, $r_2=(-1,1)$, and $(w_1,w_2)=(2,1)$.
The certificate is $2g_1+(g_2-g_1)\ge3$;
its two tight rows give $g_1=1$, $g_2=2$.

\begin{theorem}[A sparse unit-triangular certificate]\label{thm:sparse}
Under the hypotheses of Lemma~\ref{lem:prefix}, reorder the gap
coordinates so that $a_1<\cdots<a_k$. For each $i$, take the row
$h_i=e_i-v_{a_i-1}$, with $v_0=0$ as above. These are $k$ tight
chamber rows. Their matrix $H$ is unit lower triangular, with all
off-diagonal entries nonpositive. The vector
\[
 \lambda=(H^{\mathsf T})^{-1}{\bf1}
\]
has positive integer entries and satisfies
\[
 \sum_i\lambda_i h_i={\bf1},\qquad \sum_i\lambda_i=L.
\]
Thus exactly $k$ independent tight rows suffice to certify the
unique minimum, for every Singer ruler and every other complete-prefix ruler.
\end{theorem}
\begin{proof}
Adjacent gaps are distinct because all ordinary distances are distinct.
The interval of length $a_i-1$ uses only gaps smaller than $a_i$;
after the coordinate permutation its indicator therefore has support
strictly below $i$. Since $v_{a_i}=e_i$, the row $h_i$ is
$r_{a_i}$, is a chamber row, and has value one at $a$.
Write $H=I-B$ with $B$ strictly lower triangular and nonnegative
integer. Then $B^k=0$, so
\[
 (H^{\mathsf T})^{-1}=I+B^{\mathsf T}+\cdots+(B^{\mathsf T})^{k-1}.
\]
It follows that each $\lambda_i$ is an integer at least one and
$H^{\mathsf T}\lambda={\bf1}$. Since $Ha={\bf1}$,
$\sum_i\lambda_i=\lambda^{\mathsf T}Ha={\bf1}^{\mathsf T}a=L$.
Equality in the objective forces all $k$ positively weighted rows
to be tight, and invertibility gives $g=a$. Fewer than $k$ rows
cannot span the $k$-dimensional gap space.
\end{proof}

\section{Scope and further questions}

The largest-gap condition matters. Notik's twelve-mark example
\cite[Proposition~2, pp.~5--6]{notik} is a cut of a perfect difference set
modulo $133$ with reference marks
\[
 (0,4,24,27,67,77,86,88,103,119,125,132).
\]
Its removed gap is $1$, whereas its largest cyclic gap is $40$.
The real ruler
\[
 (0,4,22,25,\tfrac{123}{2},\tfrac{141}{2},\tfrac{157}{2},
 \tfrac{161}{2},94,\tfrac{217}{2},114,121)
\]
belongs to the same unit-separated chamber and is shorter.
This is Notik's prior counterexample to an arbitrary-cut extension;
that extension is not asserted in the source.

A cyclic Sidon set is a subset of $\Z/M\Z$ in which each nonzero
residue has at most one representation as an ordered difference of
distinct elements. For a prime power $q$ and a primitive
$\theta\in\mathbb F_{q^2}^{\times}$, write
\[
 B(\theta)=\{b\in\Z/(q^2-1)\Z:\theta^b=\theta+c
                  \text{ for some }c\in\mathbb F_q\}.
\]
This is Bose's cyclic Sidon construction~\cite[pp.~1--15]{bose};
see also Ruzsa~\cite[p.~267]{ruzsa} for its parameters.
For a prime $p\ge3$ and a primitive root $g$ modulo $p$, Ruzsa's
construction~\cite[Theorem~4.4 and its proof, pp.~267--268]{ruzsa}
is the cyclic Sidon set $R(p,g)\subset\Z/(p(p-1))\Z$ containing
the CRT solution $s_i\equiv i\pmod{p-1}$,
$s_i\equiv g^i\pmod p$ for each $0\le i<p-1$.
Here a Bose or Ruzsa ruler means a unit multiple of the corresponding
cyclic Sidon set, cut at a largest cyclic gap and translated to start at
zero. Notik, Section~10, page~8, asks only for
``the same for the rulers of Bose [4] and Ruzsa [15]'';
this family definition is our reading of that question.

\begin{corollary}[The Bose and Ruzsa questions]\label{cor:other}
The answer to each analogous universal chamber-uniqueness question
for the Bose and Ruzsa rulers in \cite[Section~10, p.~8]{notik} is no,
even with a largest-gap cut.
\end{corollary}
\begin{proof}
For Ruzsa, take $p=5$ and primitive root $2$. The CRT conditions
$s\equiv i\pmod4$, $s\equiv2^i\pmod5$, $0\le i<4$,
give residues $\{3,14,16,17\}\pmod{20}$. The largest gap is $11$,
from $3$ to $14$. Cutting there gives $(0,2,3,9)$, whereas
$(0,2,3,8)$ is shorter. In the reference-labelled order its distances
are $(1,2,3,5,6,8)$, compared with $(1,2,3,6,7,9)$;
both are unit separated.

For Bose, work in $\mathbb F_7[\beta]/(\beta^2-3)$ and put
$\theta=1+\beta$. The polynomial is irreducible since $3$ is a
nonsquare modulo $7$. Multiplication by $\theta$ sends $(u,v)$ to
$(u+3v,u+v)$. This recurrence gives
$\theta^{48}=1$, $\theta^{24}=6$, $\theta^{16}=4$, proving order $48$,
and the exponents of $\theta+c$, $c=0,\ldots,6$, are respectively
$(1,18,21,27,6,31,20)$. Multiplication of these exponents by the unit
$5$ gives $\{4,5,9,11,30,39,42\}\pmod{48}$.
This instance uses multiplier $5$, but
$5B(\theta)=8+B(\theta^{11})$ in $\Z/48\Z$, where
$B(\theta^{11})=\{1,3,22,31,34,44,45\}$ and $\theta^{11}$ is
primitive, so the same counterexample also belongs to the construction
without multipliers.
The largest gap is $19$, from $11$ to $30$. Its cut gives
\[
 x=(0,9,12,22,23,27,29),\qquad
 y=(0,\tfrac{17}2,\tfrac{23}2,21,22,26,28).
\]
In the labelled order determined by $x$, the distances of $y$ are
\[
\begin{split}
 &(1,2,3,4,5,6,7,\tfrac{17}2,\tfrac{19}2,\tfrac{21}2,
 \tfrac{23}2,\tfrac{25}2,\tfrac{27}2,\tfrac{29}2,\\
 &\hspace{2em}\tfrac{33}2,\tfrac{35}2,\tfrac{39}2,21,22,26,28).
\end{split}
\]
They are unit separated in exactly that order, proving $y\in C(x)$
and a strict shortening from $29$ to $28$.
The missing distances $4$ for Ruzsa and $8$ for Bose explain why
the complete-prefix argument does not apply. These feasible shorter
points are not claimed to be the chamber minima.
\end{proof}

\begin{center}
\small
\begin{tabular}{@{}>{\raggedright\arraybackslash}p{.39\linewidth}>{\raggedright\arraybackslash}p{.54\linewidth}@{}}
\toprule
Source item in \cite{notik} & Coverage here\\
\midrule
Section~7 conjecture
& Proved: Corollary~\ref{cor:singer}, for every prime power and unit multiplier.\\
Sections~7 and~10 request for structured Singer certificates
& Answered: yes. Theorem~\ref{thm:sparse} gives $k$ unit-triangular rows
and positive integer weights.\\
Section~6, Proposition~2: a shorter ruler for a non-largest cut
& Prior proof cited: Notik (2026). The arbitrary-cut extension is not
asserted in the source; it fails in this example.\\
Section~10: analogous Bose and Ruzsa questions
& Answered: no. Corollary~\ref{cor:other} supplies one largest-gap
counterexample for each family.\\
Section~7, reference~[12]: Conway--Guy and Bohman certificates
& Prior proof cited: Notik (2026), Theorems~3.2, 4.3 and~4.4 of~\cite{subset};
different subset-sum chambers.\\
Section~6 rounding remark; Section~8 asymptotic bound
& Prior proofs cited: Notik (2026), not consequences of our chamber theorem.\\
Section~10: global real/integer equality for all $n$; cases $11,12$
& Not covered; chamber uniqueness does not compare different chambers.\\
Section~10: classification of minimizers at $8,9,10$
& Not covered by this argument; no global classification is claimed.\\
Section~10: proximity of nonintegral chamber optima to the global optimum
& Not covered; the two strict shortenings give feasible chamber points,
not their optimal lengths.\\
Sections~1 (p.~1) and~5 (p.~4), reference~[13]: OEIS A003022 and A237749
& Not covered: optimal integer lengths and counts of distance-order types
are different from the minimum within a fixed chamber.\\
\bottomrule
\end{tabular}
\end{center}

\paragraph{Mathematical verification.}
The pinned Lean source~\cite{lean} contains the declaration
\begin{center}\small
\texttt{D5.S3.Geometry.FiniteGeometry.NotikSingerChamber.result}.
\end{center}
Its hypotheses are strictly increasing natural marks starting at zero,
a positive cut gap bounding every adjacent gap, and the cyclic perfect
difference property modulo the last mark plus the cut gap.
Its conclusions are uniqueness of ordinary distances, feasibility in the
real chamber, and equality of every mark whenever a chamber point has
length at most the reference length.
The finite-field specialization and the explicit dual-certificate identities
are given by the written proofs here, as are the complete-prefix
generalization and stability bounds. The second pinned Lean
source~\cite{leanother} contains
\texttt{D5.S3.Geometry.FiniteGeometry.\allowbreak{}NotikBoseRuzsaChamber.result};
it proves the negation of the disjunction of the universal Bose and
Ruzsa statements, using precisely the field and CRT witnesses in
Corollary~\ref{cor:other}.
The accompanying \texttt{verify.py} checks all perfect difference sets
containing zero of cardinalities $2$ through $6$, every largest-gap cut
of each, both certificate formulas, complete-prefix rulers with up to
five marks and length at most $18$, and the displayed counterexamples,
using exact integers, finite-field pairs and fractions.

\appendix
\section{The perfect difference property of Singer sets}

Let $F=\mathbb F_q$ and $E=\mathbb F_{q^3}$, and choose a generator
$\alpha$ of the cyclic group $E^\times$.
The quotient $E^\times/F^\times$ is cyclic of order
$N=q^2+q+1$. Let $T(z)=z+z^q+z^{q^2}$ be the trace and put
\[
 D=\{d\in\Z/N\Z:T(\alpha^d)=0\}.
\]
The condition is independent of the exponent representative, because
$\alpha^N\in F^\times$ and $T$ is $F$-linear.
The trace is not zero: its defining nonzero polynomial has degree
$q^2<q^3$, so cannot vanish on all of $E$.
Thus $H=\ker T$ is two-dimensional over $F$, and its nonzero vectors
give $(q^2-1)/(q-1)=q+1$ projective points, hence $|D|=q+1$.

For a nonzero residue $t\pmod N$, the two planes
$H$ and $\alpha^{-t}H$ are distinct. Otherwise the two nonzero linear
functionals $T(z)$ and $T(\alpha^t z)$ would have the same kernel,
so $T(\alpha^t z)=cT(z)$ for some $c\in F^\times$ and every $z$.
This would give $T((\alpha^t-c)z)=0$ for every $z$.
Since multiplication by any nonzero element permutes $E$ and $T\ne0$,
it follows that $\alpha^t=c$, contradicting $t\ne0\pmod N$.
The intersection of the two distinct planes in a three-dimensional
space is a line. Its single projective point gives exactly one residue
$d$ such that both $d$ and $d+t$ lie in $D$.
These are exactly the ordered pairs representing the residue $t$.
This establishes the perfect difference property used in
Corollary~\ref{cor:singer}, the classical construction of Singer~\cite{singer}.

\begingroup\raggedright
\begin{thebibliography}{10}
\bibitem{notik}
E.~Notik, \emph{Golomb rulers with real marks: the real and integer optima
agree up to ten marks}, Zenodo version~2 (31 October 2026),
\url{https://doi.org/10.5281/zenodo.23844781};
arXiv:2611.22656v1 (17 November 2026),
\url{https://arxiv.org/abs/2611.22656v1}.

\bibitem{bose}
R.~C.~Bose, An affine analogue of Singer's theorem,
\emph{J. Indian Math. Soc. (N.S.)} \textbf{6} (1942), 1--15.

\bibitem{ruzsa}
I.~Z.~Ruzsa, Solving a linear equation in a set of integers~I,
\emph{Acta Arith.} \textbf{65}(3) (1993), 259--282;
Theorem~4.4 and proof, pp.~267--268.
\url{https://doi.org/10.4064/aa-65-3-259-282}.

\bibitem{singer}
J.~Singer, A theorem in finite projective geometry and some applications
to number theory, \emph{Trans. Amer. Math. Soc.} \textbf{43} (1938),
377--385. \url{https://doi.org/10.1090/S0002-9947-1938-1501951-4}.

\bibitem{beck}
M.~Beck, T.~Bogart and T.~Pham, Enumeration of Golomb rulers and acyclic
orientations of mixed graphs, \emph{Electron. J. Combin.}
\textbf{19}(3) (2012), P42; arXiv:1110.6154,
\url{https://arxiv.org/abs/1110.6154}.

\bibitem{subset}
E.~Notik, \emph{Positive dual certificates for real separated subset sums:
chamber optimality of the Conway--Guy and Bohman sets, and exact values
up to eight}, arXiv:2611.17360v1 (28 October 2026),
\url{https://arxiv.org/abs/2611.17360v1}.

\bibitem{dimitro}
A.~Dimitromanolakis, \emph{Analysis of the Golomb Ruler and the Sidon
Set Problems, and Determination of Large, Near-Optimal Golomb Rulers},
diploma thesis, Technical University of Crete (June 2002),
\url{https://www.cs.toronto.edu/~apostol/golomb/main.pdf}.

\bibitem{oeisgap}
M.~Becker, \emph{Largest possible gap between adjacent elements of
cyclic planar difference sets}, OEIS A401165, comment on ruler length,
\url{https://oeis.org/A401165}.

\bibitem{lean}
\emph{Largest-gap chamber theorem}, Lean source,
declaration \texttt{NotikSingerChamber.result},
\url{https://github.com/the-omega-institute/trureturing/blob/a1ecee6d9de8eff7f246e1569bbd6211f2c471cd/D5/S3/Geometry/FiniteGeometry/NotikSingerChamber.lean}.

\bibitem{leanother}
\emph{Bose and Ruzsa chamber counterexamples}, Lean source,
declaration \texttt{NotikBoseRuzsaChamber.result},
\url{https://github.com/the-omega-institute/trureturing/blob/1a9d63013ef1b813c385f84da25351bc88cc7a98/D5/S3/Geometry/FiniteGeometry/NotikBoseRuzsaChamber.lean}.
\end{thebibliography}
\endgroup
\bigskip
\noindent\small\copyright\ 2026 Haobo Ma and Wenlin Zhang.
Licensed under Creative Commons Attribution 4.0 International
(CC BY 4.0).
\end{document}
